\section{Architecture: a two-layer pipeline}
\label{sec:architecture}

Epistract implements a six-stage pipeline (Figure~\ref{fig:architecture}). Documents enter; an agent-based extractor produces typed entities and relations against a domain schema; a deterministic validation step verifies molecular identifiers (drug-discovery only); a graph builder deduplicates and detects communities; an \emph{epistemic classifier} annotates every relation with a categorical status; and an analyst-persona stage produces a proactive briefing from the classified graph. The visualizer surfaces both layers together, with an interactive chat panel grounded in the same domain persona.

\subsection{Agent-based extraction}

Extraction agents run as Claude Code subprocesses, not stateless API calls. Each agent owns one full document end-to-end, reads it with domain understanding informed by the domain's \texttt{SKILL.md} extraction prompt, and produces a single \texttt{DocumentExtraction} JSON object conforming to the domain schema's entity and relation types. For corpora of four or more documents, agents execute in parallel, with one extractor per document. Each agent emits, per relation, a verbatim evidence quote and a confidence score in $[0, 1]$ -- two fields the epistemic classifier in Layer 2 depends on.

\paragraph{Chunking strategy.} Within each agent's lifetime, the document is read via sentence-aware chunking. The framework segments documents into chunks of up to 10{,}000 characters using \texttt{chonkie}'s \texttt{SentenceChunker}~\cite{chonkie2024} with a character tokenizer, with the final $\sim$1{,}500 characters (the last three sentences) of each chunk repeated as overlap into the next. Sentence boundaries are preserved -- chunks never split mid-sentence -- so phrases like \emph{the same compound}, \emph{this trial}, or \emph{the previous result} resolve through the overlap window rather than evaporating at chunk boundaries. The sentence-aware chunker is the v3.0 contribution to this stage of the pipeline (project planning reference: Phase~14 / FIDL-03): earlier framework versions used naive character-position splits that cut sentences and dropped cross-chunk relations; the move to \texttt{chonkie} preserves sentence integrity and adds the explicit overlap window, measurably improving cross-chunk relation recovery on long patent claims and multi-section trial protocols where coreference frequently spans paragraphs. The chunk size and overlap are tuned to fit comfortably within Claude's context window with substantial headroom for the extraction prompt and schema -- a single agent processes its document's chunks sequentially within one session, maintaining context across them.

\paragraph{Where Epistract diverges from RAG: extraction, then a second analytical pass.} The ingestion-and-chunking shape Epistract uses is not unique to it. Well-engineered RAG implementations also ingest documents one at a time, also use sentence-aware chunking with overlap (\texttt{chonkie} or an equivalent), also handle a wide range of document formats. If Epistract's contribution were simply ``we chunk well,'' there would be no contribution. The differentiation is what happens \emph{after} the chunks are read.

A RAG system embeds its chunks and indexes them; at query time, it retrieves the relevant chunks and asks an LLM to answer from the retrieved text. The product is a per-query answer; the persistent artifact is a vector index over text fragments. Epistract instead extracts \emph{typed, evidence-grounded relations} from each chunk and assembles them into a knowledge graph: the chunks are the substrate; the graph is the product. The graph carries entity types, relation types, confidence scores, verbatim evidence quotes, and -- after Layer 2 -- categorical epistemic statuses on every relation. The persistent artifact is a structured analytical record, not a vector store.

The framework then continues to analyze its own graph. The Layer 2 rule engine (Section~\ref{sec:epistemic}) classifies every relation; the cross-source aggregation step flags triples whose confidence varies across sources; the narrator stage reads the classified graph (\emph{not} the source text) and produces an analyst briefing that surfaces prophetic concentrations, contested edges, gap signals, and follow-up questions worth chasing; the workbench chat answers reactive questions against the same classified graph and the same domain persona. This second pass over the graph -- not ``what did the corpus say'' but ``what does the graph contain, and what else might it answer'' -- is the thinking process Epistract layers on top of extraction. RAG ends at retrieval. Epistract continues into analysis.

\subsection{Domain pluggability}

Domain knowledge enters the pipeline through four files, all in a self-contained per-domain directory:

\begin{itemize}
\item \texttt{domain.yaml} -- entity types and relation types in YAML, with names, descriptions, and example surface forms. The \emph{schema}.
\item \texttt{SKILL.md} -- the extraction prompt the agent reads. Includes naming conventions (INN names for drugs, HGNC symbols for genes, MeSH preferred terms for diseases), confidence-scoring guidance, and per-entity-type examples.
\item \texttt{epistemic.py} -- the per-domain rule engine that produces categorical statuses. Includes \texttt{HEDGING\_PATTERNS} (regex over evidence text), \texttt{infer\_doc\_type()} (inference rules from filename or text features), and any domain-specific custom rules (e.g., the four-level FDA epistemology classifier in \texttt{fda-product-labels/epistemic.py}).
\item \texttt{workbench/template.yaml} -- the analyst persona, in one YAML field as a multi-line string. This is the system prompt for both the reactive chat and the proactive narrator.
\end{itemize}

The core pipeline (\texttt{core/}) is domain-agnostic. \texttt{core/domain\_resolver.py} loads a domain by name, locates the four files, and threads the schema, prompt, and rule module through the rest of the pipeline. Adding a domain is a purely declarative act: write the four files, place them under \texttt{domains/<name>/}, and \texttt{/epistract:ingest --domain <name>} works. The framework ships with a \texttt{/epistract:domain} wizard that bootstraps these four files from a sample corpus through a guided three-pass interactive flow.

\subsection{Deterministic validation (drug-discovery only)}

Drug-discovery alone among the four shipped domains carries a deterministic post-extraction validation step. SMILES strings are validated and canonicalized through RDKit~\cite{rdkit2024}; protein and nucleotide sequences are checked against standard alphabets through Biopython~\cite{biopython2009}; structured identifiers (CAS registry numbers, NCT trial numbers, InChIKeys) are validated by regex and, where possible, cross-referenced. The validation report is auto-emitted as \texttt{validation\_report.json} alongside the graph. The other three shipped domains do not require this stage because their entity types are not chemistry-bound; the framework simply skips the validation step when the domain provides no \texttt{validation-scripts/} directory.

\subsection{Graph construction}

Once extraction and (optional) validation are complete, sift-kg~\cite{sifkg2024} handles graph assembly: entity deduplication via SemHash and Unicode normalization, Louvain community detection, and interactive HTML visualization via pyvis~\cite{pyvis2018}. A custom labeling engine assigns descriptive community names from member entity composition rather than numeric identifiers, with the labeling rules type-aware: gene-dominant clusters receive names reflecting the biological context of their gene members, pathway-driven clusters are named for the dominant pathway, and clusters organized around a mechanism--cell type axis receive compound labels. The output of this stage is \texttt{graph\_data.json} -- the canonical Layer 1 artifact.

\subsection{Epistemic classification (Layer 2)}

The new architectural element in v3 is a separate post-graph stage that annotates Layer 1 with categorical epistemic statuses. This is the subject of Section~\ref{sec:epistemic}; in summary, the stage reads \texttt{graph\_data.json}, applies the domain's rule engine (\texttt{epistemic.py}) to every relation's verbatim evidence and source-document type, and writes \texttt{claims\_layer.json}. This artifact is the canonical Layer 2 output. It is reproducible (rule-based) and bit-identical across runs given identical Layer 1 input.

\subsection{Persona-driven narration}

A second post-graph stage consumes the classified graph (\emph{not} the source documents) and produces a proactive briefing. The narrator's system prompt is the domain's analyst persona, loaded directly from \texttt{workbench/template.yaml}. Its input is a structured summary of the classified graph -- per-status counts, contested-claim pairs, prophetic-claim concentrations, gap signals -- assembled by \texttt{core/label\_epistemic.py} from \texttt{claims\_layer.json}. The output is \texttt{epistemic\_narrative.md}, a domain-shaped analyst briefing with executive summary, contested claims, prophetic claims, gaps, and recommended follow-ups. Unlike Layer 2, this output is non-deterministic (LLM-generated); the deterministic record is \texttt{claims\_layer.json}.

\subsection{Workbench: the reactive surface}

A web-based dashboard (the workbench, launched via \texttt{/epistract:dashboard}) renders the graph with a chat panel beside it. The chat's system prompt is the same analyst persona used for the narrator. Both surfaces -- the auto-generated briefing and the chat -- are grounded in the same graph data and the same persona. This is the persona dual-use pattern: one source of truth (\texttt{template.yaml}'s persona field), two surfaces (proactive narrator, reactive chat). Section~\ref{sec:epistemic} returns to this point.

\subsection{Plugin delivery}

Epistract is distributed as a Claude Code plugin~\cite{claudecode2024}. Users install with \texttt{/plugin install epistract} and invoke pipeline stages via slash commands: \texttt{/epistract:ingest}, \texttt{/epistract:epistemic}, \texttt{/epistract:dashboard}. The plugin encapsulates the entire dependency tree (Python virtualenv, sift-kg, RDKit, Biopython, optional sentence-transformers) and exposes a single entry point per pipeline stage.
